\documentclass[conference]{IEEEtran}
\usepackage{amsmath,amsfonts}
\usepackage{algorithmic}
\usepackage{array}
\usepackage[caption=false,font=normalsize,labelfont=sf,textfont=sf]{subfig}
\usepackage{textcomp}
\usepackage{stfloats}
\usepackage{url}
\usepackage{verbatim}
\usepackage{graphicx}
\hyphenation{op-tical net-works semi-conduc-tor IEEE-Xplore}
\def\BibTeX{{\rm B\kern-.05em{\sc i\kern-.025em b}\kern-.08em
    T\kern-.1667em\lower.7ex\hbox{E}\kern-.125emX}}
\usepackage{balance}
\begin{document}
\title{Roteamento de veículos no mercado automobilístico\\[0.5ex]
\large \textit{Um problema de logística e pesquisa operacional}}
\author{\IEEEauthorblockN{1\textsuperscript{st} Elaine Barbosa Lima}
\IEEEauthorblockA{
\textit{Universidade Federal de Minas Gerais}\\
Belo Horizonte, Brasil \\
elainebarbosalimaa@gmail.com}
\and
\IEEEauthorblockN{2\textsuperscript{nd} Erik Avelino Vaz Rocha}
\IEEEauthorblockA{
\textit{Universidade Federal de Minas Gerais}\\
Belo Horizonte, Brasil \\
@gmail.com}
\and
\IEEEauthorblockN{3\textsuperscript{rd} Lucas Henrique Martins Petronilho}
\IEEEauthorblockA{
\textit{Universidade Federal de Minas Gerais}\\
Belo Horizonte, Brasil \\
lucasgsg15@gmail.com@gmail.com}
}

\maketitle

\begin{abstract}
This document describes the most common article elements and how to use the IEEEtran class with \LaTeX \ to produce files that are suitable for submission to the Institute of Electrical and Electronics Engineers (IEEE).  IEEEtran can produce conference, journal and technical note (correspondence) papers with a suitable choice of class options.
\end{abstract}

\begin{IEEEkeywords}
Class, IEEEtran, \LaTeX, paper, style, template, typesetting.
\end{IEEEkeywords}


\section{Instrodução}

\section{The Design, Intent and \\ Limitations of the Templates}

\section{\LaTeX \ Distributions: Where to Get Them}

\section{Where to get the IEEEtran Templates}

\section{Where to get \LaTeX \ help - user groups}

\section{Document Class Options in IEEEtran}



\section{How to Create Common Front Matter}

\subsection{Author Names and Affiliations}

\subsection{Running Heads}

\subsection{Copyright Line}
\noindent For Transactions and Journals papers, this is not necessary to use at the submission stage of your paper. The IEEE production process will add the appropriate copyright line. If you are writing a conference paper, please see the ``IEEEtran\_HOWTO.pdf'' for specific information on how to code "Publication ID Marks".

\subsection{Abstracts}

\subsection{Index Terms}

\subsection{Initial Drop Cap Letter}

\subsection{Sections and Subsections}

\subsection{Citations to the Bibliography}

\subsection{Figures}

\subsection{Tables}



\begin{thebibliography}{1}

\bibitem{ams}
{\it{Mathematics into Type}}, American Mathematical Society. Online available: 

\end{thebibliography}


\end{document}


